\documentclass[11pt,a4paper]{article}

% ---------- packages ----------
\usepackage[margin=1in]{geometry}
\usepackage[T1]{fontenc}
\usepackage[utf8]{inputenc}
\usepackage{lmodern}
\usepackage{microtype}
\usepackage{amsmath,amssymb}
\usepackage{booktabs}
\usepackage{array}
\usepackage{tabularx}
\usepackage{longtable}
\usepackage{enumitem}
\usepackage{graphicx}
\usepackage{caption}
\usepackage{xcolor}
\usepackage{hyperref}
\usepackage{url}
\hypersetup{
    colorlinks=true,
    linkcolor=black,
    urlcolor=blue,
    citecolor=black,
}
\usepackage[numbers]{natbib}

% ---------- typography ----------
\setlength{\parskip}{4pt}
\setlength{\parindent}{0pt}

% Inline code
\newcommand{\code}[1]{\texttt{#1}}

% ---------- title ----------
\title{\textbf{Community-Aware Early Warning of Distrust Emergence\\
in Temporal Bitcoin Trust Networks} \\[4pt]
\large Project Milestone Report}
\author{}
\date{}

\begin{document}
\maketitle

% =====================================================================
\section{Introduction}
% =====================================================================

Online trust networks let users assign signed ratings to each other: a
positive rating expresses trust, a negative rating expresses distrust.
Bitcoin-OTC, an over-the-counter cryptocurrency marketplace, maintained such
a web-of-trust so that users could assess the reputation of potential trading
counterparts before transacting off-exchange. The resulting dataset
\citep{kumar2016} is one of the canonical benchmarks for \textbf{signed,
weighted, directed, timestamped} social network analysis.

The practical question behind this work is \textbf{early warning}: given the
state of the network at the end of a time window, can we predict where
\textit{new negative edges} will emerge in the next window? Being able to
rank-order pairs of users by their risk of entering a distrust relationship
has direct applications in marketplace moderation, fraud prevention, and
reputation governance. The problem is substantially harder than classical
link prediction because the target event is (i) \emph{rare} ($\approx 10\%$
of edges are negative overall, and only a tiny fraction of possible pairs
acquire a new negative edge per window), (ii) \emph{sign-dependent}
(predicting any new edge and predicting a new negative edge are very
different tasks), and (iii) \emph{temporally causal} (the evaluation must
respect the arrow of time).

\subsection{Research question}
\begin{quote}
At snapshot $t$, can we predict whether an unconnected or
positively-connected ordered pair $(u, v)$ will receive a \textbf{new
negative} directed edge in the next monthly window $(t, t{+}1]$, using only
information available at time $t$?
\end{quote}

\subsection{Contribution}
This project addresses the question through an \textbf{interpretable,
feature-based pipeline} aligned with the course topics --- not with heavy
deep learning. The deliverables are: (a) a temporal link-prediction setup
with leak-free cumulative monthly snapshots and a contiguous-block
train / validation / test split; (b) four compact feature families (classical
link prediction, signed structural, interpretable trust summaries, temporal /
recency, community-aware); (c) Louvain community detection on the
positive-edge trust graph with modularity diagnostics; (d) a three-way
ablation comparing structural / structural$+$temporal /
structural$+$temporal$+$community feature sets under Random Forest and
Logistic Regression; and (e) an interpretation of the weak-tie hypothesis in
a signed-network setting.

We deliberately position this work as \textbf{lighter and more interpretable}
than deep signed graph approaches such as those surveyed in
\citet{chen2025}, following course constraints against
GCN / GraphSAGE / GGNN.

% =====================================================================
\section{Related Work}
% =====================================================================

\textbf{Kumar \textit{et al.} (2016) --- edge weight prediction in weighted
signed networks \citep{kumar2016}.} Introduced two node-level trust
statistics computed iteratively on the signed graph: \textit{fairness} (how
reliable the ratings a node gives are) and \textit{goodness} (how trustworthy
the node itself is in the eyes of its raters). We borrow the \emph{spirit} of
these summaries --- a small set of interpretable per-node rating descriptors
--- rather than reimplementing the full mutually-recursive definition.
Specifically, we use \textbf{mean-rating-given},
\textbf{mean-rating-received}, and \textbf{negative-given} /
\textbf{negative-received} ratios as lightweight
``fairness-like / goodness-like'' features.

\textbf{Bertazzi \textit{et al.} (2018) --- temporal and behavioral
perspective on Bitcoin-OTC \citep{bertazzi2018}.} Characterized user-level
dynamics in Bitcoin-OTC, including how rating behavior evolves over time,
how reciprocity plays out, and how the negative-edge community grows. We use
this as motivation for the \textbf{cumulative monthly snapshot} formulation
and for including temporal features such as \code{days\_since\_last\_activity}
and exponentially decayed activity aggregates.

\textbf{Choudhury (2024) --- community-aware temporal information for
prediction \citep{choudhury2024}.} Proposed that community structure, once
extracted from the temporal graph, carries signal for link-prediction tasks
beyond what pair-level topological heuristics provide. We adopt this as the
rationale for running Louvain on every snapshot and engineering
community-aware pair features (\code{same\_community},
\code{community\_size\_u/v}, \code{boundary\_frac},
\code{neighborhood\_overlap}).

\textbf{Chen \textit{et al.} (2025) --- deep signed graph learning survey
(positioning reference) \citep{chen2025}.} Characterizes the landscape of
GNN-based signed link-prediction methods. We use this strictly for
positioning: our work is intentionally interpretable and shallow, and we do
not compare against these systems directly. Where a deep model would learn
embeddings end-to-end, we hand-engineer an analogous but human-readable
feature set.

% =====================================================================
\section{Dataset and Data-Collection Process}
% =====================================================================

\subsection{Source}
The dataset is \code{soc-sign-bitcoinotc.csv} from SNAP \citep{kumar2016}. It
was placed in \code{data/raw/} and is \emph{not modified}; all cleaning is
idempotent and produces a typed parquet file at
\code{data/interim/edges\_clean.parquet}.

\subsection{Structure (verified on disk)}
Each row is a signed rating event: \code{(source, target, rating, timestamp)}.
The file has \textbf{no header}, integer IDs, integer rating in
$[-10,+10]$, and a Unix-seconds float timestamp. Summary statistics:

\begin{center}
\begin{tabular}{lr}
\toprule
Quantity & Value \\
\midrule
Rows                         & 35{,}592 \\
Unique nodes                 & 5{,}881 \\
Positive edges (rating $> 0$)& 32{,}029 ($\approx 90\%$) \\
Negative edges (rating $< 0$)& 3{,}563 ($\approx 10\%$) \\
Zero ratings                 & 0 \\
Rating range                 & $[-10, +10]$ \\
Time span                    & 2010-11-08 to 2016-01-25 ($\approx 63$ months) \\
Reciprocity                  & 0.79 \\
Giant weakly connected component & 5{,}875 nodes \\
Giant strongly connected component & 4{,}709 nodes \\
\bottomrule
\end{tabular}
\end{center}

\subsection{Cleaning decisions}
\begin{itemize}[nosep]
  \item Parse with \code{header=None}; declare dtypes explicitly.
  \item Convert \code{timestamp} to tz-aware UTC \code{pd.Timestamp}.
  \item Defensive filters (\code{rating != 0}, \code{source != target}); none
        fire on this dataset but document the assumption.
  \item Sort by (\code{ts}, \code{source}, \code{target}) with a stable sort
        so \code{drop\_duplicates(keep="last")} produces deterministic
        ``latest rating wins''.
  \item Derive binary $\mathrm{sign} = \pm 1$.
\end{itemize}

\subsection{Cumulative monthly snapshots}
Month-end boundaries are calendar months in UTC (63 snapshots,
$t = 0, \dots, 62$). At each $t$ we build:
\begin{itemize}[nosep]
  \item $G^t_{\mathrm{dir}}$ --- directed graph of $(u, v)$ pairs with edge
        weight equal to the \textbf{most recent rating} before
        $\mathrm{end\_of\_month}(t)$ (last-rating-wins collapse, for
        snapshot graph state only).
  \item $G^t_{\mathrm{pos}}$, $G^t_{\mathrm{neg}}$ --- undirected weighted
        projections keeping only positively-signed / negatively-signed
        collapsed edges.
  \item $\mathrm{activity}[t]$ --- per-node $\{\mathrm{in},\mathrm{out},
        \mathrm{last\_ts}\}$ counters derived from the \textit{full} event
        history ($\mathrm{ts} \leq \mathrm{end\_of\_month}(t)$), not
        collapsed.
\end{itemize}
Snapshot pickles are persisted under \code{data/interim/snapshots/snap\_NNN.pkl}.
Figure~1 (\code{results/figures/edges\_per\_month.png}) shows monthly
interaction volume; Figure~2 (\code{rating\_histogram.png}) shows the rating
distribution ($+1$ and $+10$ dominate, with a long left tail).

% =====================================================================
\section{Problem Formulation}
% =====================================================================

\subsection{Core filtering}
A node $n$ is \textbf{core at snapshot $t$} iff (i) its total interactions
(in + out) up to $\mathrm{end\_of\_month}(t)$ is at least
$\mathrm{MIN\_ACTIVITY} = 3$, and (ii) its most recent activity is within
$\mathrm{RECENT\_DAYS} = 180$ of $\mathrm{end\_of\_month}(t)$. Core sizes grow
from $\approx 200$ (early) to $\approx 1{,}200$ (mid) and then contract in the
dead tail of the dataset.

\subsection{Candidate pair generation}
For each snapshot $t$, candidate ordered pairs $(u, v)$ satisfy:
\begin{enumerate}[nosep]
  \item $u, v \in \mathrm{Core}_t$, $u \neq v$;
  \item $v$ is within $\mathrm{K\_HOP} = 2$ hops of $u$ on the undirected
        any-sign projection of $G^t_{\mathrm{dir}}$;
  \item no existing \textbf{negative} directed edge $u \to v$ in
        $G^t_{\mathrm{dir}}$. Pairs with prior positive $u \to v$ are kept,
        because a rating flip from positive to negative in the next window
        legitimately counts as \emph{distrust emergence};
  \item total candidate count capped at $\mathrm{MAX\_CANDIDATES} = 200{,}000$
        via fixed-seed uniform down-sampling of negatives when exceeded.
\end{enumerate}

\subsection{Label definition}
\begin{equation*}
y(u, v) =
\begin{cases}
1 & \text{iff } \exists\ (u, v, r, \mathrm{ts}) \text{ with } r < 0 \\
  & \text{and } \mathrm{end\_of\_month}(t) < \mathrm{ts} \leq \mathrm{end\_of\_month}(t{+}1) \\
0 & \text{otherwise}.
\end{cases}
\end{equation*}
Both ``no new edge'' and ``new positive edge'' map to $y = 0$. Labels are
computed from raw events, not the collapsed snapshot.

\subsection{Train / validation / test split}
A snapshot is \textbf{eligible} iff (i) it has at least
$\mathrm{MIN\_POS\_FOR\_SPLIT} = 10$ raw new-negative events in its next
window, and (ii) it sits in the \textbf{longest contiguous block} of eligible
snapshots (this drops a sporadic late-dataset island at $t = 53$). The
resulting pool is $t \in [13, 49]$.

\begin{center}
\begin{tabular}{lcccc}
\toprule
Split & Snapshots        & Pairs       & Positives & Pos.\ rate \\
\midrule
Train & $t = 13,\dots,45$ & 5{,}013{,}346 & 545 & 0.0109\% \\
Val   & $t = 46, 47$      &   184{,}648 &  16 & 0.0087\% \\
Test  & $t = 48, 49$      &   148{,}639 &  34 & 0.0229\% \\
\bottomrule
\end{tabular}
\end{center}
The split is \textbf{locked} from Stage~2 onward; no tuning or feature design
ever sees test.

\subsection{Leakage controls}
\begin{itemize}[nosep]
  \item All features at snapshot $t$ are computed from
        $\mathrm{df\_up\_to\_t} = \{\text{rows with }
        \mathrm{ts} \leq \mathrm{end\_of\_month}(t)\}$. A defensive
        \code{assert df\_up\_to\_t["ts"].max() <= end\_ts} is embedded in
        both the temporal and trust-summary precompute functions and is
        re-checked at every snapshot by the Stage 3--5 drivers.
  \item Labels are derived strictly from events in
        $(\mathrm{end\_of\_month}(t), \mathrm{end\_of\_month}(t{+}1)]$.
  \item Splits are block-temporal: no overlap by construction. An
        end-to-end sanity check prints
        \code{[OK] no train/val/test snapshot overlap}.
\end{itemize}

\subsection{Class-imbalance handling}
Training negatives are subsampled to a \textbf{1:100 positive:negative ratio}
(keeping all 545 positives), shrinking the training matrix from ${\sim}5\mathrm{M}$
to ${\sim}55\mathrm{k}$ rows; validation and test remain full. Ranking metrics
(ROC-AUC, PR-AUC, Precision@k) are preferred over threshold-dependent
precision/recall.

% =====================================================================
\section{Methodology}
% =====================================================================

\subsection{Feature families}
\begin{center}
\small
\begin{tabular}{lrll}
\toprule
Family & \# feats & Source & Status \\
\midrule
Classical link prediction    & 4 & undirected any-sign projection        & implemented \\
Signed structural            & 8 & directed graph + $G_{\mathrm{pos}}/G_{\mathrm{neg}}$ & implemented \\
Interpretable trust summary  & 4 & raw events (fairness/goodness-like)   & implemented \\
Temporal / recency           & 12 & raw events, recency windows, decay   & implemented \\
Community-aware              & 8 & Louvain partition of $G_{\mathrm{pos}}$ & implemented \\
Iterated fairness/goodness   & 2 & mutually-recursive update             & \textbf{deferred} \\
\bottomrule
\end{tabular}
\end{center}
The full feature list (36 columns) is canonicalized in
\code{src/build\_features.py :: FEATURE\_COLUMNS}.

\subsection{Models}
Three models are trained per feature-set ablation:
\begin{itemize}[nosep]
  \item \textbf{Random baseline} (uniform probabilities, seeded) --- calibrates
        metric floors.
  \item \textbf{Logistic Regression} --- \code{StandardScaler} $+$
        \code{class\_weight='balanced'}, \code{liblinear} solver, 2000 iters.
  \item \textbf{Random Forest} --- 200 trees,
        \code{min\_samples\_leaf = 20},
        \code{class\_weight = 'balanced\_subsample'}.
\end{itemize}
No hyperparameter tuning beyond these defaults. This is intentional: the
point of this project is feature engineering and ablation, not model
selection.

\subsection{Evaluation protocol}
Primary metrics: \textbf{ROC-AUC}, \textbf{PR-AUC} ($=$ average precision),
\textbf{Precision@k} for $k \in \{50, 100\}$. Threshold-0.5
precision / recall / F1 are reported for completeness only; they are not
informative at $0.01\%$ base rate. The primary operating mode is ranking.

% =====================================================================
\section{Mathematical Background}
% =====================================================================

\subsection{Classical link-prediction heuristics}
Let $N(x)$ denote the neighborhood of $x$ in the undirected projection at
snapshot $t$.
\begin{align*}
\mathrm{CN}(u, v) &= |N(u) \cap N(v)| \\
J(u, v)           &= \frac{|N(u) \cap N(v)|}{|N(u) \cup N(v)|} \\
\mathrm{AA}(u, v) &= \sum_{w \in N(u) \cap N(v)} \frac{1}{\log |N(w)|} \\
\mathrm{PA}(u, v) &= |N(u)| \cdot |N(v)|
\end{align*}
Adamic-Adar \citep{adamic2003} is a weighted CN that discounts hub neighbors;
PA models ``rich get richer'' dynamics.

\subsection{Louvain modularity maximization}
Louvain \citep{blondel2008} greedily optimizes weighted modularity
\begin{equation*}
Q = \frac{1}{2m} \sum_{i, j} \left[ A_{ij}
   - \frac{k_i k_j}{2m} \right] \delta(c_i, c_j),
\end{equation*}
where $A_{ij}$ is the edge weight between nodes $i, j$, $k_i$ the weighted
degree of $i$, $m = \tfrac{1}{2} \sum_{i,j} A_{ij}$, and $c_i$ the community
of $i$. We apply Louvain on the per-snapshot \textbf{positive-edge
undirected weighted graph} $G^t_{\mathrm{pos}}$ (weights $=$ count of directed
positive edges collapsed into the undirected pair). Random seed is fixed;
outputs are cached per snapshot in \code{data/interim/communities/}.

Across the 37 split snapshots, modularity is consistently in
$[0.47, 0.52]$, with 15--47 communities and top-5 sizes ranging from
${\sim}300$ to ${\sim}1{,}200$ nodes --- confirming that $G_{\mathrm{pos}}$ has
well-defined community structure throughout the study period.

\subsection{Evaluation metrics}
\begin{itemize}[nosep]
  \item \textbf{ROC-AUC}: area under the ROC curve; probability that a random
        positive is ranked above a random negative.
  \item \textbf{PR-AUC} ($=$ average precision): area under the
        precision-recall curve; at extreme imbalance, this is a stricter
        summary than ROC-AUC.
  \item \textbf{Precision@k}: fraction of true positives in the top $k$
        scored candidates. A ranking-oriented metric that maps directly onto
        the ``watchlist of $k$ risky pairs'' early-warning use case.
\end{itemize}

\subsection{Exponential decay weighting}
For each event at time $\mathrm{ts}$ and snapshot boundary
$\mathrm{end\_ts}$, define
$\Delta = (\mathrm{end\_ts} - \mathrm{ts}).\mathrm{days}$ and weight
$w = \exp(-\Delta / \tau)$. Per-node decayed aggregates (total activity,
negative count) are then $\sum_e w_e$ grouped by source or target. We use
$\tau = 60$ days (a Bitcoin-OTC-appropriate half-life given the burstiness of
user activity).

\subsection{Neighborhood overlap (weak-tie diagnostic)}
\begin{equation*}
\mathrm{NO}(u, v) = \frac{|N(u) \cap N(v) \setminus \{u, v\}|}
                         {|N(u) \cup N(v) \setminus \{u, v\}|}.
\end{equation*}
Granovetter-style \citep{granovetter1973}: a weak tie bridges otherwise
weakly-connected neighborhoods and therefore has \emph{low} overlap. We
expect this feature to rank highly in a community-aware feature set if the
weak-tie hypothesis applies to distrust emergence.

% =====================================================================
\section{Preliminary Findings and Summary Statistics}
% =====================================================================

\subsection{Network-level statistics (whole period)}
Reported in \S3.2. The graph is large, mostly connected (giant WCC is $99.9\%$
of nodes), and strongly reciprocal ($r = 0.79$). The positive-edge subgraph
is dense; the negative-edge subgraph is sparse.

\subsection{Per-snapshot summary (Stage 2)}
Persisted at \code{results/stage2\_summary.csv}. Key observations:
\begin{itemize}[nosep]
  \item \textbf{Core size} ramps from dozens ($t = 0$) to ${\sim}1{,}200$
        ($t \approx 28$) and then contracts as the dataset ages out.
  \item \textbf{Candidate count} is dominated by 2-hop reachability on the
        core and hits the 200{,}000 cap for roughly half the mid-period
        snapshots.
  \item \textbf{Raw next-window negative count} (our eligibility proxy)
        fluctuates between 10--250 during the active period and collapses to
        0--9 past September 2015 --- the reason the split pool is truncated
        to $t \in [13, 49]$.
  \item \textbf{Core filtering} is the dominant positive-loss step (e.g., at
        $t = 21$, 244 raw new-negatives $\to 54$ in-core $\to 41$
        in-candidates). The 2-hop restriction loses only 0--10 positives per
        snapshot and is \emph{not} the bottleneck.
\end{itemize}

\subsection{Louvain community snapshot (Stage 5)}
\begin{center}
\small
\begin{tabular}{rrrrl}
\toprule
$t$ & $|V(G_{\mathrm{pos}})|$ & \# comm. & modularity & top-5 sizes \\
\midrule
13 & 1{,}631 & 19 & 0.488 & $321/212/202/176/125$ \\
28 & 3{,}652 & 19 & 0.484 & $665/601/510/488/343$ \\
44 & 5{,}317 & 39 & 0.513 & $994/956/870/763/654$ \\
49 & 5{,}455 & 47 & 0.515 & $1{,}107/1{,}030/962/679/572$ \\
\bottomrule
\end{tabular}
\end{center}
Modularity is stable in $[0.47, 0.52]$ and community count grows gracefully
with the network.

\subsection{Positive rate: same- vs cross-community (TRAIN)}
\begin{center}
\begin{tabular}{lrrr}
\toprule
Partition                & Pairs       & Positives & Pos.\ rate \\
\midrule
\code{same\_community = 1} & 1{,}540{,}969 & 215 & \textbf{0.0140\%} \\
\code{same\_community = 0} & 3{,}472{,}377 & 330 & 0.0095\% \\
\bottomrule
\end{tabular}
\end{center}
\textbf{Same-community pairs show a $\approx 47\%$ higher positive rate} than
cross-community pairs --- the \emph{opposite} of the naive Granovetter
prediction at community granularity. We elaborate in \S7.5.

\subsection{Early ablation (validation; not final test numbers)}
These are reported to show the feature families are behaving as expected.
The final numbers will be reported on the held-out test split.

\begin{center}
\small
\begin{tabular}{lccccc}
\toprule
Model & Struct. & $+$ Temp. & $+$ Comm. & $\Delta(T-S)$ & $\Delta(C-T)$ \\
\midrule
LogReg (ROC-AUC)        & 0.897  & \textbf{0.944} & 0.942 & $+0.047$   & $-0.002$        \\
LogReg (PR-AUC)         & 0.0032 & 0.0022         & 0.0021 & ---        & ---             \\
Random Forest (ROC-AUC) & 0.868  & 0.955          & \textbf{0.958} & $+0.087$   & $+0.003$    \\
Random Forest (PR-AUC)  & 0.0017 & 0.0029         & 0.0034 & $+71\%$ rel. & $+17\%$ rel. \\
\bottomrule
\end{tabular}
\end{center}

\textbf{Temporal features drive the main lift}, as expected from
Bertazzi-style behavioral dynamics. \textbf{Community features deliver a
smaller incremental gain}, concentrated in PR-AUC --- the metric that
matters under extreme imbalance. The most useful \emph{community-aware}
feature is \code{neighborhood\_overlap}, whose LR coefficient is the
strongest inside the community family. This provides \textbf{partial
support} for the weak-tie interpretation: at the pair level, Granovetter's
pattern holds (lower shared-neighborhood $=$ higher risk of new negative
edge), but at the community level the intuitive direction reverses (distrust
surfaces \emph{inside} a cluster, not at its boundary). The synthesis we
adopt for the final narrative is \textbf{``intra-community weak ties''} ---
pairs with few shared neighbors nested inside the same trusted cluster.

\subsection{Feature separation (TRAIN, positives vs negatives)}
Top positive-vs-negative mean gaps across families:
\begin{itemize}[nosep]
  \item \textbf{Structural:} \code{pa} (4{,}185 vs 915), \code{pos\_out\_u}
        (51.9 vs 25.5), \code{neg\_out\_u} (14.1 vs 2.8),
        \code{reciprocity\_flag} (0.30 vs 0.04).
  \item \textbf{Trust summary:} \code{mean\_rating\_received\_v} (0.17 vs
        1.05 --- victim-side signal is the strongest interpretable trust
        descriptor).
  \item \textbf{Temporal:} \code{days\_since\_last\_activity\_u} (15 vs 62
        --- recency is the single strongest temporal signal),
        \code{recent\_total\_given\_u} (17 vs 4).
  \item \textbf{Community:} \code{neighborhood\_overlap} (0.044 vs 0.059),
        \code{boundary\_frac\_u} (0.41 vs 0.36).
\end{itemize}

% =====================================================================
\section{General Difficulties}
% =====================================================================

\begin{enumerate}[label=\arabic*.,leftmargin=*]
  \item \textbf{Class imbalance.} Base rate is $\approx 0.01\%$ on the
        candidate set. Most standard metrics (accuracy, threshold-based
        precision/recall/F1) are uninformative. Mitigation:
        training-negative subsampling (1:100) plus ranking metrics
        (ROC-AUC, PR-AUC, Precision@k).
  \item \textbf{Tail of the dataset is dead.} Bitcoin-OTC activity collapses
        after mid-2015, leaving the last ${\sim}12$ months with 0--9 new
        negatives per window. A na\"ive ``last two snapshots'' test fold
        yields 0 positives in validation. We address this with an
        eligibility rule ($\geq 10$ raw new-negatives in next window) and a
        \textbf{longest-contiguous-block} split-pool constraint. The
        eligible pool ends at $t = 49$ (April 2015).
  \item \textbf{Candidate explosion.} All-pairs predictions are
        $O(n^2) \approx 3.5 \times 10^7$ per snapshot. Mitigation: core
        filtering (active nodes only) $+$ 2-hop restriction $+$ $200\mathrm{k}$
        cap with fixed-seed downsampling.
  \item \textbf{Leakage risk.} Any feature touching post-$t$ events would
        corrupt the evaluation. Mitigation: a single-choke-point
        \code{df\_up\_to\_t} argument passed into every feature precompute,
        defensive asserts, and a per-snapshot post-hoc
        $\mathrm{max\_ts\_used} \leq \mathrm{end\_ts}$ check.
  \item \textbf{Multi-edges and rating flips.} On Bitcoin-OTC, a pair can be
        re-rated. For the \textbf{snapshot graph state} we apply
        \emph{last-rating-wins} (current trust judgment). For
        \textbf{temporal / trust history features} we use the full event
        history. These two conventions must be kept consistent across
        features --- they are, by design separation of the \code{snap}
        pickle from the raw \code{df\_up\_to\_t} slice.
  \item \textbf{Interpretability budget.} We committed to avoiding
        GCN / GraphSAGE / deep signed models. This keeps the story coherent
        and the features legible at the cost of absolute predictive power.
        Results so far suggest the interpretable setup is competitive enough
        for the early-warning use case.
  \item \textbf{Louvain determinism.} Louvain is stochastic. We fix
        \code{random\_state} and cache partitions per snapshot; all runs
        from Stage~5 onward reference the cached partition.
\end{enumerate}

% =====================================================================
\section{What Remains and Plan to Completion}
% =====================================================================

\subsection{Done so far (core pipeline)}
\begin{itemize}[nosep]
  \item Stage 1 --- data loading, cleaning, EDA, snapshot generation.
  \item Stage 2 --- core filtering, candidate generation, label generation,
        split selection, sanity checks.
  \item Stage 3 --- classical link-prediction $+$ signed structural $+$
        trust-summary features; first baselines.
  \item Stage 4 --- temporal features; feature correlation / redundancy
        analysis; $1{:}100$ negative-subsampling training protocol.
  \item Stage 5 --- Louvain community detection with cached metadata;
        community-aware features; 3-way validation ablation.
  \item Engineering: deterministic seeds, cached intermediate artifacts,
        leakage sanity checks at every stage.
\end{itemize}

\subsection{Remaining for the final report}
\begin{enumerate}[nosep,label=\arabic*.]
  \item \textbf{Held-out test evaluation.} Final ablation
        (structural / $+$temporal / $+$temporal$+$community) on the test
        split ($t = 48, 49$). Early indicative numbers exist but will be
        finalized into the report-ready comparison table.
  \item \textbf{Error analysis.} Top-ranked true-positive and false-positive
        pairs, with a short discussion of the patterns (recurring ``victim''
        nodes, recurring ``issuer'' nodes, cross- vs same-community,
        overlap distributions).
  \item \textbf{Feature importance summary}, grouped by family (structural /
        temporal / community), for both LR (standardized coefficients) and
        RF (Gini importance). Used both for interpretation and for
        identifying redundant features (e.g., \code{recent\_*} vs
        \code{decayed\_*} pairs are $\geq 0.9$ correlated --- we will note
        this openly rather than prune after the fact).
  \item \textbf{Report figures.} Finalize: (i) degree distribution on log
        axes, (ii) modularity-over-time curve, (iii) core-size /
        candidate-count / positives-count per-snapshot panel, (iv) ROC and
        PR curves for the three feature sets on test, (v) a same-community
        vs cross-community positive-rate bar chart.
  \item \textbf{Report writing.} Introduction, related work, methodology,
        results, discussion, limitations, conclusion. Target length: 8--10
        pages (IEEE style).
  \item \textbf{Slide deck} for the presentation, drawn from the final
        report.
\end{enumerate}

\subsection{Optional (only if time permits)}
\begin{itemize}[nosep]
  \item \textbf{node2vec $+$ Logistic Regression} as a learned-embedding
        baseline, strictly as an optional extension gated behind a working
        main pipeline.
  \item \textbf{Motif / triad counts} (signed triangles in the vicinity of
        the candidate pair).
  \item \textbf{Bitcoin-Alpha transfer test}: fit on Bitcoin-OTC and score
        on the sibling Bitcoin-Alpha dataset as a zero-shot generalization
        check.
  \item \textbf{Iterated Kumar-style fairness/goodness} (2--3 iterations,
        explicitly not to convergence). Deferred: the four simple
        trust-summary features already cover the interpretable slot; a
        richer iterative version would be a 2-feature add-on if the final
        numbers motivate it.
\end{itemize}

\subsection{Explicitly out of scope}
GNN / GraphSAGE / GCN / GGNN / other heavy deep learning methods. Null-model
comparisons (considered at the outset but not planned for the final report).

% =====================================================================
\section{Repository and Reproducibility}
% =====================================================================

Project tree (abbreviated):
\begin{verbatim}
462 project/
|-- data/ {raw, interim, processed}/
|-- src/
|   |-- config.py, io_utils.py, cleaning.py, eda.py
|   |-- snapshots.py, core.py, candidates.py, labels.py
|   |-- features_structural.py, features_trust.py
|   |-- features_temporal.py
|   |-- communities.py, features_community.py
|   |-- build_features.py, models.py
|-- scripts/ run_stage{1..7}_*.py
|-- results/ {figures, stage2_summary.csv, ...}
|-- report/ milestone_report.tex, (final_report.*)
|-- requirements.txt, README.md, .gitignore
\end{verbatim}

All intermediate artifacts are regenerable from the raw CSV via the
\code{run\_stage\{1..7\}\_*.py} scripts; random seeds are fixed in
\code{src/config.py :: RANDOM\_SEED = 42}.

% =====================================================================
\begin{thebibliography}{9}
\bibitem{kumar2016}
S.~Kumar, F.~Spezzano, V.~S. Subrahmanian, and C.~Faloutsos,
\emph{Edge weight prediction in weighted signed networks},
IEEE ICDM, 2016.

\bibitem{bertazzi2018}
I.~Bertazzi \textit{et al.},
\emph{Temporal and behavioral perspective on Bitcoin-OTC}, 2018.

\bibitem{choudhury2024}
N.~Choudhury,
\emph{Community-aware temporal network prediction}, 2024.

\bibitem{chen2025}
Z.~Chen \textit{et al.},
\emph{Deep signed graph learning: a survey and benchmark}, 2025.

\bibitem{granovetter1973}
M.~Granovetter,
\emph{The strength of weak ties},
American Journal of Sociology, 1973.

\bibitem{blondel2008}
V.~D. Blondel, J.-L. Guillaume, R.~Lambiotte, and E.~Lefebvre,
\emph{Fast unfolding of communities in large networks},
J.\ Stat.\ Mech., 2008.

\bibitem{adamic2003}
L.~A. Adamic and E.~Adar,
\emph{Friends and neighbors on the web},
Social Networks, 2003.

\bibitem{easleykleinberg}
D.~Easley and J.~Kleinberg,
\emph{Networks, Crowds, and Markets},
Cambridge University Press, 2010 (Chs. 3, 5).
\end{thebibliography}

\textit{Bibliographic formatting will be standardized in the final report.}

\end{document}
